\scanpage{35}
\deleted{On suppose la trace continue.}
\margin{\uline{Astuce}: on \unsure{plonge} un \unsure{anneau} avec une trace \unsure{non} continue, dans un \unsure{anneau} avec une trace continue.}
\margin{\deleted{$\unsure{1+}\varphi_{\unsure{|A|}}(\unsure{t})$}}

\uline{Prop. 1}
\[
\varphi_A(\unsure{t+0})=
\inf_{\substack{E\in P\\\operatorname{Tr}E<t}}
\sup_{x\in(1-E)H}(Ax,x)
\qquad(\text{trace continue})
\]
On a en effet, comme on voit aussitôt, (décomposition spectrale)
$\varphi_A(t)\geq$ \unsure{second membre}, il suffit donc de prouver
\[
\sup_{x\in(1-E)H}(Ax,x)\geq\varphi_A^-(t)
\qquad\inserted{\text{si }\operatorname{Tr}E<t}.
\]
On est ramené au cas où $A$ est égal à $\lambda F$, $F$ projecteur de trace $\geq t$, donc à \unsure{trace} $=$.
\inserted{(car $A\geq\varphi_A^{\deleted{\unsure{+}}\unsure{-}}(t)F$ avec $\operatorname{Tr}F\geq t$)}
\deleted{\unsure{On voit}}
Il suffit de prouver: si \uline{$\operatorname{Tr}E<\operatorname{Tr}F$, alors $(1-E)\mathcal H\wedge F\mathcal H\ne(0)$}.

En effet, on a \deleted{$\unsure{E}$} pour
\[
\mathcal X=E\mathcal H,\quad\mathcal Y=F\mathcal H,\qquad
\mathcal X'=\mathcal H\ominus\mathcal X,\quad
\mathcal Y'=\unsure{\mathcal H|\mathcal X}\ominus\mathcal Y:
\]
\[
\mathcal X\ominus(\mathcal X\wedge\mathcal Y')\sim
\mathcal Y\ominus(\mathcal Y\wedge\mathcal X')
\]
donc $\mathcal Y\wedge\mathcal X'=0$ impliquerait
$\deleted{\mathcal X\ominus}\mathcal Y\unsure{\prec|\lesssim}\mathcal X$
d'où $\operatorname{Tr}F\leq\operatorname{Tr}E$, absurde.

\uline{Corollaire 1} $A$ quelconque \unsure{implique}:
\[
\varphi_{|A|}(t+0)=\inf_{\substack{E\in P\\\operatorname{Tr}E<t}}
\sup_{x\in(1-E)\mathcal H}\|Ax\|
\qquad(\text{trace continue})
\]

\uline{Corollaire 2}
\[
\begin{aligned}
\varphi_{|PA|}\deleted{\unsure{(t)}}&\leq\|P\|\varphi_{|A|}\deleted{\unsure{(t)}}
\qquad(\text{trace quelconque}),\\
\varphi_{|AQ|}&\leq\|Q\|\varphi_{|A|}
\end{aligned}
\]
\margin{\deleted{$\unsure{2|x}\unsure{\sqrt{A}}\ \unsure{A\in\underline b}$}}

\uline{Corollaire 3} $\varphi_A$ croît avec l'hermitien \unsure{positif} $A$.
\margin{$\varphi_{\alpha(A)}=\alpha\varphi_A$}
\margin{\uline{Corollaire 4}
$\varphi_{|1+A|}\leq\varphi_{1+|A|}=1+\varphi_{|A|}$:
\unsure{preuve}
\deleted{$\Delta_t(1+A)\leq\Delta_t(1+|A|)$}
\inserted{$\unsure{\varphi|\Phi}_{|1+A|}\leq\unsure{\varphi|\Phi}_{\unsure{1+|A|}}$}}
\diagram{\begin{tikzcd}[row sep=small,column sep=small]
&\text{Corollaire 3}\arrow[d,bend right=70]\\
&\text{Corollaire 5}\\
{}\arrow[ur]&
\end{tikzcd}}

\uline{Corollaire 5} \unsure{Équicontinuité} \inserted{des \unsure{fonctions} $A\to\varphi_{|A|}(t)$}
\deleted{$\unsure{A\to}\varphi_{|A|}(t+0)$ \unsure{fonction de} $A$.}

\uline{Prop. 2} $\varphi_{|A|}=\varphi_{|A^*|}$.
Il suffit de \unsure{le} prouver \unsure{que} $\varphi_{AA^*}=\varphi_{A^*A}$
\unsure{de le} prouver pour $A\in\unsure{\underline a|a}$.
\unsure{Prenons}
\deleted{$A\in\underline a$, \unsure{et l'on a} \unsure{pour}}
$\unsure{\mu}_{A^*A}=\unsure{\mu}_{AA^*}$
\deleted{$\unsure{=\int x^p\,d\mu_{|A|}(x^p)}$}
\mspage{36}
